\part{Operational Features}

\chapter{Signal K and Data Health}
\label{ch:datahealth}

At sea, a number is only useful if you know whether it is still fresh.
FairWindSK uses three data states consistently throughout the interface
and operational documentation.

\section{Individual Value States}

\rowcolors{2}{LtGray}{white}
\begin{longtable}{L{2.2cm}L{5.4cm}L{6.1cm}}
\toprule
\rowcolor{Navy}
\color{white}\sffamily\bfseries State &
\color{white}\sffamily\bfseries Meaning &
\color{white}\sffamily\bfseries Recommended action \\
\midrule
\textbf{Live}    &
  Value received recently and updating normally. &
  Use normally. Continue to glance at the Top Bar periodically. \\
\textbf{Stale}   &
  Updates have slowed or stopped. Last known value shown, dimmed.
  A warning, not current truth. &
  Treat with caution. Cross-check with independent instruments.
  Investigate if the state persists. \\
\textbf{Missing} &
  No usable value. A dash or placeholder is shown.
  \textbf{Missing is not zero.} It means unavailable. &
  Treat as entirely unknown. Never use as a safe default.
  Verify with another source. \\
\bottomrule
\end{longtable}

\section{Connection Health States}

\rowcolors{2}{LtGray}{white}
\begin{longtable}{L{3.2cm}L{4.8cm}L{5.7cm}}
\toprule
\rowcolor{Navy}
\color{white}\sffamily\bfseries State &
\color{white}\sffamily\bfseries Meaning &
\color{white}\sffamily\bfseries Response \\
\midrule
Live               & Fully healthy. &
  Continue normally. \\
Connecting         & Establishing the initial connection. &
  Wait 30--60 seconds. \\
Reconnecting       & Connection lost; auto-recovery in progress. &
  Wait briefly. Usually resolves automatically. \\
Stale              & Connected but updates slowed or delayed. &
  Cross-check instruments. Investigate if it persists. \\
REST degraded      & Request/response channel impaired. &
  Settings pages may fail; monitor closely. \\
Stream degraded    & Live stream impaired. &
  Readouts may stop updating; treat as stale. \\
Disconnected       & No connection. &
  Check Signal~K server power and Wi-Fi. \\
Apps loading       & App catalogue being fetched. &
  Wait; apps appear when done. \\
Apps stale         & App catalogue may be out of date. &
  Apps still work; list may be incomplete. \\
Foreground degraded& Current application has a problem. &
  Retry once; if still failing, return to launcher and check shell health. \\
\bottomrule
\end{longtable}

\section{When to Act Immediately}

Act without delay when:

\begin{itemize}
  \item depth becomes stale or missing in shallow or restricted water
  \item position or heading disappears during close-quarters navigation
  \item a POB event occurs
  \item an autopilot command does not match the vessel's actual response
\end{itemize}

Monitor briefly and then cross-check when:

\begin{itemize}
  \item the shell is reconnecting after a brief drop
  \item app loading takes longer than expected
  \item route guidance values stop updating
  \item only one app looks degraded while the rest of the shell is healthy
\end{itemize}

\section{Automatic Recovery Cycle}

When a connection is lost, FairWindSK automatically:

\begin{enumerate}
  \item Detects the disconnection or degradation.
  \item Performs a fresh Signal~K server discovery (handles endpoint changes after a restart).
  \item Re-establishes the WebSocket stream.
  \item Re-sends all active subscriptions.
  \item Performs an immediate REST snapshot so widgets show current values right away.
  \item Refreshes the app catalogue and all My~Data resource models.
  \item Refreshes unit preferences from the server.
\end{enumerate}

When a Signal~K server restart is triggered from inside FairWindSK (via the Signal~K
admin web app), FairWindSK detects the planned restart, pauses its reconnect loop
briefly, and then resumes recovery automatically.

After any reconnect or server restart, wait for the recovery to complete and confirm
that live data returns before depending on safety-critical values.

\begin{warningbox}
After a power interruption, GPS receivers and instruments take time to re-acquire and
settle.
Even after FairWindSK shows \textit{Live}, verify each safety-critical value before
acting on it.
\end{warningbox}

\chapter{Person Over Board (POB)}
\label{ch:pob}

\begin{warningbox}
Mark a person overboard event immediately.
Do not delay because you expect to remember the position later.
Mark first, then continue recovery actions.
The POB function assists with recovery reference.
It does not replace immediate crew response, lookout, helmsmanship, or established
emergency procedures.
\end{warningbox}

\section{Activating POB}

Tap \iconlabel{alarm-pob}{POB} in the Bottom Bar.
It is permanently visible and never disabled.
Practise finding it in the dark and in rough conditions.

In the instant you tap POB, FairWindSK simultaneously:

\begin{itemize}
  \item Captures the current GPS position as the incident position.
  \item Creates a timestamped waypoint resource on Signal~K named
        \code{POB HH:MM:SS} (UTC).
  \item Sets that waypoint as the active navigation destination ---
        all navigation apps immediately show bearing and distance to the POB.
  \item Raises the standard Signal~K \skpath{notifications.mob} alarm at emergency
        level, alerting all connected devices and displays.
  \item Starts the elapsed-time clock counting up every second.
  \item Subscribes to live bearing and distance from the vessel to the POB waypoint.
  \item Replaces the normal Bottom Bar with the POB dialog.
\end{itemize}

\screenshot{pob-bar-active}{POB Bar: bearing, distance, elapsed time, SAR search buttons}{}

\section{POB Bar Controls and Readouts}

\rowcolors{2}{LtGray}{white}
\begin{longtable}{C{1.4cm}L{3.0cm}L{9.3cm}}
\toprule
\rowcolor{Navy}
\color{white}\sffamily\bfseries Icon &
\color{white}\sffamily\bfseries Control &
\color{white}\sffamily\bfseries Description \\
\midrule
---                      & POB selector     &
  Drop-down listing all active POB events.
  Select an event to display its data. \\
---                      & Last known position &
  Latitude and longitude at the moment POB was pressed,
  in the configured coordinate format. \\
---                      & Bearing          &
  Live compass bearing from vessel to POB waypoint.
  Updates continuously as the vessel manoeuvres. \\
---                      & Distance         &
  Live distance from vessel to POB waypoint. Updates continuously. \\
---                      & Elapsed time     &
  HH:MM:SS counting up from the moment POB was pressed. \\
---                      & Cancel           &
  Cancels the selected incident.
  Removes the Signal~K alarm and navigation destination for that event. \\
\iconlg{search-spiral}   & Spiral search    &
  Generates a SAR spiral route centred on the POB position. \\
\iconlg{search-parallel} & Parallel search  &
  Generates a SAR parallel-track route centred on the POB position. \\
\iconlg{close-google}    & Close dialog     &
  Returns to the normal Bottom Bar.
  \textbf{The alarm and waypoint remain active on Signal~K.} \\
\bottomrule
\end{longtable}

\begin{warningbox}
Tapping \iconlg{close-google}~Close does \textbf{not} cancel the POB.
The Signal~K alarm and waypoint remain active until explicitly cancelled.
Only the \key{Cancel} button removes a specific incident.
\end{warningbox}

\section{SAR Search Patterns}

\subsection{Spiral Search}

Tap \iconlabel{search-spiral}{Spiral}.
FairWindSK creates a $500 \times 500$\,m SAR area region centred on the POB
position and a spiral route starting at the POB and radiating outward in 50\,m
spacing legs.
Both the region and route are saved as Signal~K resource objects and appear
immediately in \ui{My~Data~> Regions} and \ui{My~Data~> Routes}.

\subsection{Parallel-Track Search}

Tap \iconlabel{search-parallel}{Parallel}.
FairWindSK creates the same $500 \times 500$\,m area region and a
creeping-line parallel-track route covering it in 60\,m-spaced parallel legs.

\section{Cancelling a POB}

\begin{enumerate}
  \item Select the incident in the POB selector (if multiple incidents are active).
  \item Tap \key{Cancel}.
  \item FairWindSK removes the Signal~K alarm and clears the navigation destination.
  \item If no more active incidents remain, the POB dialog closes automatically.
\end{enumerate}

\begin{notebox}
The POB waypoint is \textbf{not} deleted when cancelled.
It is retained in \ui{My~Data~> Waypoints} with its timestamp and position for
accident reporting and post-incident debrief.
\end{notebox}

\chapter{Autopilot Controls}
\label{ch:autopilot}

Tap \iconlabel{command-autopilot}{Autopilot} in the Bottom Bar to open the Autopilot
bar, which replaces the normal Bottom Bar.
The Autopilot icon is active only when the required autopilot plugin is installed and
FairWindSK can interact with it.

\begin{warningbox}
Always verify the real-world effect of an autopilot command on the vessel.
Never assume a mode change took effect because you tapped a control.
Be prepared to disengage or override immediately if the response is unsafe or
unexpected.
\end{warningbox}

\screenshot{autopilot-bar}{Autopilot Bar: mode buttons, trim controls, rudder indicator}{}

\section{Autopilot Bar Controls}

Icons are the exact assets defined in \code{AutopilotBar.cpp}.

\rowcolors{2}{LtGray}{white}
\begin{longtable}{C{1.4cm}L{2.4cm}L{4.6cm}L{5.3cm}}
\toprule
\rowcolor{Navy}
\color{white}\sffamily\bfseries Icon &
\color{white}\sffamily\bfseries Control &
\color{white}\sffamily\bfseries Signal~K action &
\color{white}\sffamily\bfseries When to use \\
\midrule
\iconlg{close-google}                & Stand By    & Disengage         & Any time you want manual control. Verify the helm responds. \\
\iconlg{command-autopilot}           & Auto        & Compass mode      & Open-water passages; holds current heading. \\
\iconlg{up-iec}                      & Wind        & Wind-vane mode    & Upwind/downwind; holds wind angle. \\
\iconlg{navigation-route}            & Route       & GPS route mode    & Following a waypoint route. \\
\iconlg{spinbox-minus}               & $-1\textdegree$ port  & Trim $-1\textdegree$ & Fine adjustment to port. \\
\iconlg{spinbox-plus}                & $+1\textdegree$ stbd  & Trim $+1\textdegree$ & Fine adjustment to starboard. \\
\iconlg{chevron-double-left-google}  & $-10\textdegree$ port & Adjust $-10\textdegree$ & Larger correction to port. \\
\iconlg{chevron-double-right-google} & $+10\textdegree$ stbd & Adjust $+10\textdegree$ & Larger correction to starboard. \\
\iconlg{arrow-left-google}           & Tack port   & Tack to port      & Sailing upwind. \\
\iconlg{arrow-right-google}          & Tack stbd   & Tack to starboard & Sailing upwind. \\
\iconlg{chevron-double-left-google}  & Gybe port   & Gybe to port      & Light--moderate conditions; monitor. \\
\iconlg{chevron-double-right-google} & Gybe stbd   & Gybe to starboard & Supervise at helm in strong winds. \\
\iconlg{waypoint-optional-iec}       & Next WPT    & Advance route leg & Route navigation. \\
\iconlg{alerts}                      & Dodge       & Obstacle avoidance & Lobster pot, vessel, hazard. \\
\bottomrule
\end{longtable}

\section{Rudder Indicator}

The Autopilot bar shows a \textbf{read-only} horizontal slider displaying current
rudder deflection from $-30\textdegree$ (full port) to $+30\textdegree$ (full
starboard).
This reflects the actual rudder position reported by the pilot ---
not a commanded heading.
Use it to confirm the autopilot is responding to commands as expected.

\chapter{Anchor Watch}
\label{ch:anchor}

Tap \iconlabel{anchor-iec}{Anchor} to open the Anchor Watch bar.

\begin{warningbox}
Electronic anchor watch is an aid only.
In an exposed anchorage, deteriorating weather, or poor holding ground,
post a human watch regardless of what any electronic device shows.
\end{warningbox}

\screenshot{anchor-bar}{Anchor Watch Bar: position, depth, rode, radius, windlass controls}{}

\section{Anchor Bar Controls and Readouts}

Icons are the exact assets defined in \code{AnchorBar.ui}.

\rowcolors{2}{LtGray}{white}
\begin{longtable}{C{1.4cm}L{3.0cm}L{9.3cm}}
\toprule
\rowcolor{Navy}
\color{white}\sffamily\bfseries Icon &
\color{white}\sffamily\bfseries Control &
\color{white}\sffamily\bfseries Description \\
\midrule
\iconlg{anchor-drop}    & Drop        & Records GPS position as anchor point and starts monitoring. \\
\iconlg{anchor-raise}   & Raise       & Stops monitoring and clears the anchor position. \\
\iconlg{anchor-release} & Release     & Issues a release-rode command to an electric windlass. \\
\iconlg{anchor-up}      & Up (hold)   & Continuous haul-in command to windlass. Release to stop. \\
\iconlg{anchor-down}    & Down (hold) & Continuous let-out command to windlass. Release to stop. \\
\iconlg{anchor-reset}   & Reset       & Resets anchor position to current GPS position. \\
---                      & Radius slider & Maximum distance permitted before alarm. \\
---                      & Set Radius  & Commits the radius and sends it to the anchor plugin. \\
---                      & LAT / LON   & GPS coordinates of the recorded anchor position. \\
---                      & Depth       & Current water depth from the depth sounder. \\
---                      & Rode        & Length of rode deployed. \\
---                      & Current radius / bearing & Live distance and bearing from vessel to anchor. \\
---                      & Max radius  & Maximum distance reached since watch started. \\
\bottomrule
\end{longtable}

\section{Setting Anchor Watch --- Step by Step}

\begin{enumerate}
  \item Anchor normally and pay out your intended scope.
  \item After the vessel has settled, tap \iconlg{anchor-iec} in the Bottom Bar.
  \item Confirm the displayed GPS position is where the anchor is lying.
  \item Tap \iconlg{anchor-drop}~\key{Drop}. FairWindSK records the position and starts monitoring.
  \item Set the radius, allowing for scope, swing circle, GPS accuracy ($\pm3$--$10$\,m),
        and a safety margin.
  \item Tap \key{Set Radius} to commit.
  \item Close the bar. Monitoring continues in the background.
\end{enumerate}

\chapter{Alarms}
\label{ch:alarms}

Tap \iconlabel{alarm-general}{Alarms} to open the Alarms bar.
When any alarm is active the icon is highlighted with accent colouring across all
comfort presets --- designed to catch the eye even during the night watch.

\screenshot{alarms-bar}{Alarms Bar: SOLAS emergency buttons and active alarm list}{}

\section{Emergency Alarm Buttons}

Icons are the exact assets defined in \code{AlarmsBar.ui}.

\rowcolors{2}{LtGray}{white}
\begin{longtable}{C{1.4cm}L{1.8cm}L{3.8cm}L{6.7cm}}
\toprule
\rowcolor{Navy}
\color{white}\sffamily\bfseries Icon &
\color{white}\sffamily\bfseries Button &
\color{white}\sffamily\bfseries Signal~K path &
\color{white}\sffamily\bfseries Emergency type \\
\midrule
\iconlg{alarm-pob}     & POB     & \skpath{notifications.mob}     & Person / man overboard \\
\iconlg{alarm-sinking} & Sinking & \skpath{notifications.sinking} & Vessel taking water or sinking \\
\iconlg{alarm-fire}    & Fire    & \skpath{notifications.fire}    & Fire aboard \\
\iconlg{alarm-piracy}  & Piracy  & \skpath{notifications.piracy}  & Piracy or armed attack \\
\iconlg{alarm-abandon} & Abandon & \skpath{notifications.abandon} & Abandon ship \\
\iconlg{alarm-adrift}  & Adrift  & \skpath{notifications.adrift}  & Vessel adrift and uncontrolled \\
\bottomrule
\end{longtable}

\begin{warningbox}
Activating an alarm in FairWindSK raises a Signal~K notification on the onboard
network only.
It does \textbf{not} contact coast guard, MRCC, or any external authority.
Make a DSC and/or voice distress call on VHF Channel~16 in any real emergency.
\end{warningbox}

\section{Alarm Priority and Response}

Handle alerts in this order:

\begin{enumerate}
  \item Safety-critical values first --- depth, position, heading.
  \item Connection and shell health second.
  \item App-specific problems third.
\end{enumerate}

\section{Active Alarm List}

Below the buttons, the Alarms bar lists all currently active Signal~K notifications,
showing path, state (emergency/alarm/warn), message text, and timestamp.
It also shows the current Signal~K connection state and the last stream update
timestamp --- a compact health summary without leaving the bar.
